\documentclass[../../../../master.tex]{subfiles}

\begin{document}

\monographleaf{Neural Architectures, Representations \& Regularization}
  {MLPs, Kolmogorov-Arnold Networks, Residuals \& Regularization}
  {Feedforward Networks, Normalization, Residuals \& KANs}
  {Feedforward Networks, Universal Approximation \& Kolmogorov-Arnold Networks (KAN)}
  {Prince Ch.~3--4; Bishop Ch.~6; Goodfellow Ch.~6; Math Behind KAN; Murphy Intro Sec.~13.2}

% =============================================================================
\section{Problem Formulation \& Unified Notation}
\label{sec:mlpkan:formulation}
% =============================================================================

Supervised learning seeks to approximate an unknown target mapping $f^* : \mathcal{X} \subseteq \R^{D_i} \to \R^{D_o}$ from finite observations $\{(\vecb{x}_n, \vecb{y}_n)\}_{n=1}^N$. Classical linear basis-function models parameterize predictions as $y(\vecb{x}, \vecb{w}) = \sum_{j=1}^M w_j \varphi_j(\vecb{x}) + w_0$ using \emph{fixed}, data-independent basis functions $\varphi_j$ \bookcite{deep_learning_foundations_concepts}{Sec.~6.1, p.~172}. For multivariate polynomials of degree $M$ in $D_i$ inputs, the number of independent coefficients grows as $O(D_i^M)$, triggering Bellman's curse of dimensionality \bookcite{deep_learning_foundations_concepts}{Sec.~6.1.1, p.~173}. Feedforward neural networks---specifically Multilayer Perceptrons (MLPs)---overcome this bottleneck by making the basis functions themselves \emph{data-adaptive} through compositions of affine maps and elementwise nonlinearities. Kolmogorov-Arnold Networks (KANs) offer a dual paradigm: instead of learning linear weights on fixed node activations, KANs place learnable univariate spline functions directly on the graph edges.

Let an $L$-layer architecture have layer widths $[n_0, n_1, \dots, n_L]$ with $n_0 = D_i$ and $n_L = D_o$, and set $\vecb{h}^{(0)} = \vecb{x}$.
\begin{itemize}
  \item \textbf{Multilayer Perceptron (MLP):} Each layer $\ell \in \{1,\dots,L-1\}$ applies an affine transformation followed by a fixed scalar activation $\sigma : \R \to \R$ evaluated elementwise:
  \begin{equation}
    \label{eq:mlpkan:mlp_layer}
    \vecb{a}^{(\ell)} = \mat{W}^{(\ell)} \vecb{h}^{(\ell-1)} + \vecb{b}^{(\ell)}, \qquad
    \vecb{h}^{(\ell)} = \sigma\!\left(\vecb{a}^{(\ell)}\right),
  \end{equation}
  where $\mat{W}^{(\ell)} \in \R^{n_\ell \times n_{\ell-1}}$ and $\vecb{b}^{(\ell)} \in \R^{n_\ell}$, culminating in a linear (or task-specific) readout $\hat{\vecb{y}} = \mat{W}^{(L)} \vecb{h}^{(L-1)} + \vecb{b}^{(L)}$ \bookcite{probabilisticml_kevinmurphy_introduction}{Sec.~13.2.2, p.~428}, \bookcite{understandingdl_simonprince}{Sec.~4.4.1, p.~49}.
  \item \textbf{Kolmogorov-Arnold Network (KAN):} Each layer $\ell \in \{1,\dots,L\}$ replaces scalar edge weights with a matrix of learnable univariate functions $\mat{\Phi}^{(\ell)} = [\phi_{\ell,q,p}]_{q=1..n_\ell,\, p=1..n_{\ell-1}}$, and each node performs pure summation:
  \begin{equation}
    \label{eq:mlpkan:kan_layer}
    h^{(\ell)}_q = \sum_{p=1}^{n_{\ell-1}} \phi_{\ell,q,p}\!\left(h^{(\ell-1)}_p\right), \qquad q = 1, \dots, n_\ell,
  \end{equation}
  where each $\phi_{\ell,q,p}$ is parameterized as a residual SiLU term plus a degree-$k$ B-spline on $G$ grid intervals \bookcite{math_behind_kan}{PDF p.~6}.
\end{itemize}

\begin{crossbook}[Cross-Book Notation \& Pedagogical Reconciliation]
\label{tab:mlpkan:notation}
\small
\begin{tabularx}{\textwidth}{@{} >{\raggedright\arraybackslash}p{2.7cm} >{\raggedright\arraybackslash}p{2.9cm} >{\raggedright\arraybackslash}X >{\raggedright\arraybackslash}X @{}}
\toprule
\textbf{Concept} & \textbf{Unified Monograph} & \textbf{Prince / Bishop / Goodfellow} & \textbf{Murphy / KAN Monograph} \\
\midrule
Input / Output dim. & $D_i = n_0,\; D_o = n_L$ & Prince: $D_i, D_o$; Goodfellow: $d$ inputs; Bishop: $D$ inputs & Murphy: $D$ inputs; KAN: $n_0, n_L$ \\
Layer weights & $\mat{W}^{(\ell)} \in \R^{n_\ell \times n_{\ell-1}}$ & Prince: $\mat{\Omega}_{k-1}$; Bishop/Goodfellow: $\mat{W}^{(\ell)}$ & Murphy: $\mat{W}_l$; KAN: $c_m$ (spline coeffs) \\
Layer bias & $\vecb{b}^{(\ell)} \in \R^{n_\ell}$ & Prince: $\vecb{\beta}_{k-1}$; Goodfellow: $\vecb{c}, \vecb{b}$ & Murphy: $\vecb{b}_l$; KAN: absorbed in $\phi_{\ell,q,p}$ \\
Pre- / Post-activation & $\vecb{a}^{(\ell)},\; \vecb{h}^{(\ell)}$ & Prince: $\vecb{h}_k$; Bishop: $a, z$; Goodfellow: $\vecb{h}$ & Murphy: $\vecb{a}_l, \vecb{z}_l$; KAN: $x_p, \Phi_q$ \\
Activation function & $\sigma(\cdot)$ (MLP), $\phi_{\ell,q,p}(\cdot)$ (KAN) & Prince: $a[\cdot]$; Bishop: $h(\cdot)$; Goodfellow: $g(\cdot)$ & Murphy: $\varphi_l(\cdot)$; KAN: $w_b b(x) + w_s \mathrm{spline}(x)$ \\
Depth / Width & $L$ layers, $n_\ell$ units & Prince: $K$ hidden layers, $D$ units; Goodfellow: depth $l$, width $n$ & Murphy: $L$ layers, $K_l$ units; KAN: $[n_0,\dots,n_L]$ \\
\bottomrule
\end{tabularx}
\end{crossbook}

% =============================================================================
\section{Core Mathematical Derivations \& Theorems}
\label{sec:mlpkan:theory}
% =============================================================================

\subsection{Universal Approximation Theorems for Shallow Networks}

As established in \bookcite{deep_learning_goodfellow_bengio_courville_2017}{Sec.~6.4.1, p.~194} and \bookcite{probabilisticml_kevinmurphy_introduction}{Sec.~13.2.5, p.~434}, a single hidden layer suffices for uniform approximation on compact domains whenever the activation function is non-polynomial.

\begin{coremath}[Cybenko--Hornik Universal Approximation Theorem (UAT)]
\label{thm:mlpkan:uat}
Let $K \subset \R^{D_i}$ be compact and let $f \in C(K, \R)$ be any continuous target function. Suppose $\sigma : \R \to \R$ is continuous and non-polynomial (e.g., logistic sigmoid, $\tanh$, or $\ReLU$). Then for every tolerance $\epsilon > 0$, there exist a finite width $N \in \N$, output weights $\alpha_j \in \R$, input weights $\vecb{w}_j \in \R^{D_i}$, and biases $\theta_j \in \R$ ($j = 1, \dots, N$) such that
\begin{equation}
  \label{eq:mlpkan:uat}
  \sup_{\vecb{x} \in K} \abs[\Bigg]{f(\vecb{x}) - \sum_{j=1}^N \alpha_j \sigma\!\left(\vecb{w}_j\T \vecb{x} + \theta_j\right)} < \epsilon.
\end{equation}
\end{coremath}

\begin{intuition}[Geometric Partitioning vs.\ the Width Curse]
In 1D with $\ReLU(z) = \max(0, z)$, each hidden unit $h_j(x) = \ReLU(\theta_{j0} + \theta_{j1}x)$ introduces a single ``joint'' (hinge) at $x = -\theta_{j0}/\theta_{j1}$, so $D$ hidden units partition $\R$ into at most $D + 1$ linear segments \bookcite{understandingdl_simonprince}{Sec.~3.2, p.~29}. Adding units refines the piecewise-linear interpolant to arbitrary precision. In $D_i$ dimensions, however, approximating a generic Lipschitz function to accuracy $\epsilon$ with a shallow network can require $N = \Omega(\epsilon^{-D_i})$ hidden units---an exponential bottleneck unless the target exhibits compositional structure.
\end{intuition}

\subsection{Piecewise-Linear Polyhedral Regions \& Exponential Depth Separation}

A ReLU network composes continuous piecewise-linear (CPWL) maps, so the full network $f(\vecb{x}; \vecb{\theta})$ is itself CPWL over convex polyhedral cells in $\R^{D_i}$, each corresponding to a fixed binary activation pattern $\vecb{s} \in \{0,1\}^{\sum_\ell n_\ell}$.

\begin{coremath}[Linear Region Counting \& Depth Separation]
\label{thm:mlpkan:regions}
\begin{enumerate}
  \item \textbf{1D Shallow vs.\ Deep Exact Bounds \bookcite{understandingdl_simonprince}{Sec.~4.5.2, p.~50}:} A scalar-input ($D_i = 1$) shallow ReLU network with $D > 2$ hidden units has $3D + 1$ parameters and creates at most $D + 1$ linear regions. A deep ReLU network with $K$ hidden layers of width $D > 2$ has $3D + 1 + (K-1)D(D+1)$ parameters and can create up to
  \begin{equation}
    \label{eq:mlpkan:1d_regions}
    R_{\max}(1, K, D) = (D + 1)^K \quad \text{linear regions.}
  \end{equation}
  For example, $K=5, D=10$ uses only $471$ parameters yet produces up to $11^5 = 161{,}051$ linear regions.
  \item \textbf{Multivariate Depth Separation (Mont\'ufar et al.) \bookcite{deep_learning_goodfellow_bengio_courville_2017}{Sec.~6.4.1, Eq.~6.42, p.~200}:} For a ReLU network with $d$ inputs, depth $l$, and $n \ge d$ units per hidden layer, the maximum number of linear regions scales as
  \begin{equation}
    \label{eq:mlpkan:montufar}
    R_{\max}(d, l, n) = \Omega\!\left(\left(\frac{n}{d}\right)^{(l-1)d} n^d\right),
  \end{equation}
  growing exponentially in depth $l$, whereas a shallow ($l=1$) network is bounded by Zaslavsky's hyperplane arrangement count $\sum_{j=0}^{d} \binom{n}{j} = O(n^d)$.
\end{enumerate}
\end{coremath}

\begin{intuition}[Input-Space Folding]
Each hidden layer in a deep ReLU (or absolute-value) network folds the input space along the hyperplanes $\vecb{w}_j\T \vecb{h} + b_j = 0$, mapping distinct regions onto overlapping images \bookcite{deep_learning_goodfellow_bengio_courville_2017}{Fig.~6.5, p.~200}, \bookcite{understandingdl_simonprince}{Fig.~4.3, p.~44}. Subsequent layers then slice all folded sheets simultaneously, multiplying the region count across layers rather than adding regions linearly.
\end{intuition}

\subsection{Smooth \& Piecewise-Linear Activation Taxonomy}

As noted in \bookcite{probabilisticml_kevinmurphy_introduction}{Sec.~13.2.3, p.~428}, if $\sigma$ is linear then Eq.~\eqref{eq:mlpkan:mlp_layer} collapses to a single affine map $\mat{W}' \vecb{x} + \vecb{b}'$. Classical sigmoids ($\sigma(a) = (1+e^{-a})^{-1}$ and $\tanh(a)$) saturate exponentially as $\abs{a} \to \infty$, causing vanishing gradients. Modern non-saturating activations---detailed in \bookcite{deep_learning_foundations_concepts}{Sec.~6.2.3, pp.~184--185} and \bookcite{understandingdl_simonprince}{Fig.~3.13, p.~37}---include:
\begin{align}
  \mathrm{ReLU}(a) &= \max(0, a), &
  \mathrm{LeakyReLU}_\alpha(a) &= \max(0, a) + \alpha \min(0, a), \label{eq:mlpkan:relu_leaky} \\
  \mathrm{Softplus}(a) &= \ln(1 + e^a), &
  \mathrm{GELU}(a) &= a\,\Phi(a), \label{eq:mlpkan:softplus_gelu} \\
  \mathrm{SiLU}_\beta(a) &= \frac{a}{1 + e^{-\beta a}}, &
  \sigma(a) &= \frac{1}{1 + e^{-a}}, \label{eq:mlpkan:smooth_acts}
\end{align}
where $\Phi(a)$ is the standard Gaussian CDF and $\beta=1$ gives the standard SiLU/Swish activation used inside KAN residual branches. Unlike ReLU, GELU, SiLU, and Softplus are $C^\infty$ smooth and non-monotonic near $a \in [-2, 0)$ (for GELU/SiLU), improving gradient flow for negative pre-activations.

\subsection{Kolmogorov--Arnold Superposition \& B-Spline KANs}

Rather than approximating functions via sums of ridge functions $\sigma(\vecb{w}_j\T \vecb{x} + \theta_j)$, Kolmogorov and Arnold proved an \emph{exact} decomposition into univariate functions \bookcite{math_behind_kan}{PDF p.~2}.

\begin{coremath}[Kolmogorov--Arnold Representation Theorem (1957)]
\label{thm:mlpkan:ka_theorem}
Every continuous multivariate function $f : [0,1]^n \to \R$ can be represented exactly as a two-layer superposition of univariate continuous functions:
\begin{equation}
  \label{eq:mlpkan:ka_rep}
  f(x_1, \dots, x_n) = \sum_{q=0}^{2n} \Phi_q\!\left(\sum_{p=1}^n \phi_{q,p}(x_p)\right),
\end{equation}
where $\phi_{q,p} : [0,1] \to \R$ and $\Phi_q : \R \to \R$ are continuous univariate functions.
\end{coremath}

Classical attempts to apply Eq.~\eqref{eq:mlpkan:ka_rep} directly failed because the inner functions $\phi_{q,p}$ can be fractal and non-differentiable, rendering gradient-based optimization intractable \bookcite{math_behind_kan}{PDF p.~2}. Modern KANs resolve this by (i) generalizing the fixed $[n, 2n+1, 1]$ depth-2 template to arbitrary layer widths $[n_0, n_1, \dots, n_L]$, and (ii) parameterizing each edge $\phi_{\ell,q,p}$ as a smooth combination of a residual basis $b(x) = \mathrm{SiLU}(x)$ and a degree-$k$ B-spline \bookcite{math_behind_kan}{PDF p.~6}:
\begin{equation}
  \label{eq:mlpkan:kan_edge}
  \phi(x) = w_b \, b(x) + w_s \, \mathrm{spline}(x), \qquad
  \mathrm{spline}(x) = \sum_{m=0}^{G+k-1} c_m B_{m,k}(x),
\end{equation}
where $w_b, w_s$ control the global residual magnitude and $\{c_m\}$ are learnable spline control coefficients. Given an augmented non-decreasing knot vector $t_0 \le t_1 \le \dots \le t_{G+2k}$ spanning $[a, b]$ with $G$ interior intervals, the B-spline basis functions $B_{m,k}(x)$ are defined recursively via the \textbf{Cox--de Boor formula}:
\begin{align}
  B_{m,0}(x) &= \mathbb{I}[t_m \le x < t_{m+1}], \label{eq:mlpkan:cdb0} \\
  B_{m,r}(x) &= \frac{x - t_m}{t_{m+r} - t_m} B_{m,r-1}(x) + \frac{t_{m+r+1} - x}{t_{m+r+1} - t_{m+1}} B_{m+1,r-1}(x), \quad r = 1, \dots, k. \label{eq:mlpkan:cdbr}
\end{align}
Crucially, $B_{m,k}(x)$ has \emph{local support} $[t_m, t_{m+k+1})$ (at most $k+1$ basis functions are nonzero at any $x$), is $C^{k-1}$ smooth on uniform knots, and satisfies the partition-of-unity identity $\sum_{m=0}^{G+k-1} B_{m,k}(x) = 1$ for $x \in [t_k, t_{G+k}]$.

% =============================================================================
\section{Algorithmic Implementation \& Comparative Synthesis}
\label{sec:mlpkan:algo}
% =============================================================================

Unlike MLPs---which require retraining from scratch when increasing hidden width---KANs support \textbf{Grid Extension} \bookcite{math_behind_kan}{PDF pp.~7--8}: a model trained on a coarse grid $G_1$ can be refined to a finer grid $G_2 > G_1$ by solving a local least-squares projection $\min_{\vecb{c}'} \E_{x}\!\left[\left(\sum_j c'_j B_{j,k}^{\text{fine}}(x) - \sum_m c_m B_{m,k}^{\text{coarse}}(x)\right)^2\right]$ on each edge, preserving the learned function while unlocking higher spatial resolution for fine-tuning. Algorithm~\ref{alg:mlpkan:kan_forward} details the forward evaluation and grid-extension step.

\begin{algorithm}[htbp]
\caption{KAN Layer Forward Pass (Cox--de Boor) \& Grid Extension}
\label{alg:mlpkan:kan_forward}
\KwIn{Input $\vecb{h}^{(\ell-1)} \in \R^{n_{\ell-1}}$, knot vector $\vecb{t} \in \R^{G+2k+1}$, spline coefficients $\mat{C} \in \R^{n_\ell \times n_{\ell-1} \times (G+k)}$, residual weights $\mat{W}_b, \mat{W}_s \in \R^{n_\ell \times n_{\ell-1}}$, spline degree $k$.}
\KwOut{Layer output $\vecb{h}^{(\ell)} \in \R^{n_\ell}$.}
\BlankLine
\For{$p \leftarrow 1$ \KwTo $n_{\ell-1}$}{
  $x \leftarrow h^{(\ell-1)}_p$\;
  Initialize $B_{m,0}(x) \leftarrow \mathbb{I}[t_m \le x < t_{m+1}]$ for $m = 0, \dots, G+2k-1$\;
  \For{$r \leftarrow 1$ \KwTo $k$}{
    \For{$m \leftarrow 0$ \KwTo $G+2k-1-r$}{
      Compute $B_{m,r}(x)$ via Cox--de Boor recursion in Eq.~\eqref{eq:mlpkan:cdbr}\;
    }
  }
  $b_p \leftarrow x / (1 + e^{-x})$ \tcp*{SiLU residual basis}
}
\For{$q \leftarrow 1$ \KwTo $n_\ell$}{
  $h^{(\ell)}_q \leftarrow \sum_{p=1}^{n_{\ell-1}} \left[ W_{b,q,p}\, b_p + W_{s,q,p} \sum_{m=0}^{G+k-1} C_{q,p,m}\, B_{m,k}\!\left(h^{(\ell-1)}_p\right) \right]$\;
}
\BlankLine
\textbf{Optional Grid Extension ($G \to G_{\text{new}}$):} Construct fine knots $\vecb{t}'$, sample collocation points $\{x_s\}_{s=1}^S \subset [t_k, t_{G+k}]$, and solve $\mat{C}'_{q,p} \leftarrow \argmin_{\vecb{c}'} \sum_{s=1}^S \left(\sum_j c'_j B^{\text{fine}}_{j,k}(x_s) - \sum_m C_{q,p,m} B_{m,k}(x_s)\right)^2$ via QR/normal equations\;
\end{algorithm}

\begin{table}[htbp]
\centering
\caption{Per-layer parameter and computational complexity comparison between Shallow MLPs, Deep MLPs, and B-spline KANs ($N = n_{\ell-1} = n_\ell$, spline degree $k$, grid intervals $G$).}
\label{tab:mlpkan:complexity}
\small
\begin{tabularx}{\textwidth}{@{}l l l >{\raggedright\arraybackslash}X@{}}
\toprule
\textbf{Architecture} & \textbf{Params / Layer} & \textbf{FLOPs / Sample} & \textbf{Hardware \& Inductive Trade-off} \\
\midrule
Shallow MLP ($L=2$) & $N(D_i + D_o + 1)$ & $O(N(D_i + D_o))$ & Dense GEMM; $O(\epsilon^{-D_i})$ width curse \\
Deep MLP ($L$ layers) & $N^2 + N$ & $O(N^2)$ & Dense GEMM; $O((N/D_i)^{L D_i})$ regions \\
KAN ($L$ layers, $G, k$) & $N^2(G + k + 2)$ & $O(N^2 k^2)$ (local) & Symbolic interpretability \& grid extension; non-GEMM overhead \\
\bottomrule
\end{tabularx}
\end{table}

\begin{crossbook}[Cross-Book Synthesis: Compositional Priors vs.\ Univariate Edge Plasticity]
\label{box:mlpkan:synthesis}
\bookcite{deep_learning_foundations_concepts}{Sec.~6.3.1, p.~187} and \bookcite{deep_learning_goodfellow_bengio_courville_2017}{Sec.~6.4.1, p.~201} frame deep MLPs as encoding a \emph{compositional program prior}: intermediate hidden layers extract reusable motifs (such as solving XOR via AND/OR sub-circuits \bookcite{probabilisticml_kevinmurphy_introduction}{Sec.~13.2.1, p.~427}) that exponentially multiply linear regions per parameter \bookcite{understandingdl_simonprince}{Sec.~4.5.2, p.~50}. By contrast, \bookcite{math_behind_kan}{PDF pp.~6--7} shifts representational capacity from depth/width scaling to \emph{edge-level B-spline plasticity}. When the target admits low-dimensional compositional or physical-law structure, KANs achieve lower approximation error with far fewer nodes and allow post-hoc symbolic regression on individual edges; when high-dimensional unstructured data (e.g., images or token embeddings) dominate, dense GEMM-friendly MLPs remain orders of magnitude faster per FLOP.
\end{crossbook}

% =============================================================================
\section{Visual Diagnostics \& Geometric Structure}
\label{sec:mlpkan:visuals}
% =============================================================================

Figure~\ref{fig:mlpkan:mlp_vs_kan_tikz} contrasts the computational graphs of an MLP and a KAN, highlighting the duality between fixed node nonlinearities with scalar edge weights and linear node summations with spline-parameterized edges.

\begin{figure}[htbp]
  \centering
  \begin{tikzpicture}[>=Stealth, scale=0.92, transform shape]
    % --- Left Panel: MLP ---
    \node[font=\small\bfseries\color{monographNavy}] at (2.2, 4.1) {MLP: Fixed Node Activation};
    \node[monoBox, minimum width=0.9cm, minimum height=0.6cm] (mx1) at (0.3, 2.7) {$x_1$};
    \node[monoBox, minimum width=0.9cm, minimum height=0.6cm] (mx2) at (0.3, 1.1) {$x_2$};

    \node[monoTealBox, minimum width=1.25cm, minimum height=0.75cm] (mh1) at (2.4, 3.1) {$\sigma(\cdot)$};
    \node[monoTealBox, minimum width=1.25cm, minimum height=0.75cm] (mh2) at (2.4, 1.9) {$\sigma(\cdot)$};
    \node[monoTealBox, minimum width=1.25cm, minimum height=0.75cm] (mh3) at (2.4, 0.7) {$\sigma(\cdot)$};

    \node[monoIndigoBox, minimum width=1.0cm, minimum height=0.65cm] (my) at (4.4, 1.9) {$\hat{y}$};

    \foreach \s in {mx1, mx2} {
      \foreach \t in {mh1, mh2, mh3} {
        \draw[monoArrow] (\s.east) -- (\t.west);
      }
    }
    \foreach \s in {mh1, mh2, mh3} {
      \draw[monoArrow] (\s.east) -- (my.west);
    }
    \node[font=\scriptsize\color{monographSlate}] at (1.35, 3.35) {$w_{qp}^{(1)}$};
    \node[font=\scriptsize\color{monographSlate}] at (3.45, 2.75) {$w_{1q}^{(2)}$};
    \node[font=\scriptsize\color{monographSlate}, align=center] at (2.35, -0.1) {$\vecb{h}^{(\ell)} = \sigma(\mat{W}^{(\ell)}\vecb{h}^{(\ell-1)} + \vecb{b}^{(\ell)})$};

    % Divider
    \draw[dashed, color=monographRule, line width=0.8pt] (5.55, -0.3) -- (5.55, 4.3);

    % --- Right Panel: KAN ---
    \node[font=\small\bfseries\color{monographNavy}] at (9.1, 4.1) {KAN: Learnable Edge B-Splines};
    \node[monoBox, minimum width=0.9cm, minimum height=0.6cm] (kx1) at (6.4, 2.7) {$x_1$};
    \node[monoBox, minimum width=0.9cm, minimum height=0.6cm] (kx2) at (6.4, 1.1) {$x_2$};

    \node[monoIndigoBox, minimum width=1.05cm, minimum height=0.55cm] (ks1) at (9.1, 3.0) {$\sum$};
    \node[monoIndigoBox, minimum width=1.05cm, minimum height=0.55cm] (ks2) at (9.1, 0.8) {$\sum$};

    \node[monoTealBox, minimum width=0.95cm, minimum height=0.6cm] (ky) at (11.7, 1.9) {$\sum \to \hat{y}$};

    \draw[monoArrow, color=monographTeal] (kx1.east) -- node[above, sloped, font=\scriptsize] {$\phi_{1,1}(x_1)$} (ks1.west);
    \draw[monoArrow, color=monographTeal] (kx1.east) -- (ks2.west);
    \draw[monoArrow, color=monographTeal] (kx2.east) -- (ks1.west);
    \draw[monoArrow, color=monographTeal] (kx2.east) -- node[below, sloped, font=\scriptsize] {$\phi_{2,2}(x_2)$} (ks2.west);

    \draw[monoArrow, color=monographIndigo] (ks1.east) -- node[above, sloped, font=\scriptsize] {$\Phi_1(\cdot)$} (ky.west);
    \draw[monoArrow, color=monographIndigo] (ks2.east) -- node[below, sloped, font=\scriptsize] {$\Phi_2(\cdot)$} (ky.west);

    \node[font=\scriptsize\color{monographSlate}, align=center] at (9.1, -0.1) {$\phi(x) = w_b\,\mathrm{SiLU}(x) + w_s \sum_m c_m B_{m,k}(x)$};
  \end{tikzpicture}
  \caption{Architectural duality between an MLP (fixed scalar node activations $\sigma$, linear scalar edge weights) and a KAN (linear summation at nodes, learnable univariate B-spline activations $\phi_{q,p}$ on edges).}
  \label{fig:mlpkan:mlp_vs_kan_tikz}
\end{figure}

Figure~\ref{fig:mlpkan:activations_bsplines} plots the canonical activation functions alongside cubic B-spline basis functions $B_{m,3}(x)$ on uniform knots, illustrating both the smooth non-monotonicity of GELU/SiLU and the local support and partition-of-unity property of B-splines.

\begin{figure}[htbp]
  \centering
  \begin{minipage}[t]{0.48\textwidth}
    \centering
    \begin{tikzpicture}
      \begin{axis}[
        width=\linewidth, height=5.4cm,
        xlabel={$a$}, ylabel={$\sigma(a)$},
        xmin=-3.5, xmax=3.5, ymin=-0.6, ymax=3.5,
        grid=both, grid style={monographRule, opacity=0.4},
        legend pos=north west, legend style={font=\scriptsize, draw=none, fill=white, fill opacity=0.85},
        tick label style={font=\scriptsize}, label style={font=\footnotesize}
      ]
        \addplot[domain=-3.5:3.5, samples=120, monographNavy, thick] {max(0, x)};
        \addlegendentry{ReLU}
        \addplot[domain=-3.5:3.5, samples=120, monographTeal, thick, dashed] {ln(1 + exp(x))};
        \addlegendentry{Softplus}
        \addplot[domain=-3.5:3.5, samples=120, monographIndigo, thick] {x / (1 + exp(-x))};
        \addlegendentry{SiLU}
        \addplot[domain=-3.5:3.5, samples=120, monographCrimson, thick, dotted] {0.5 * x * (1 + tanh(0.79788456 * (x + 0.044715 * x^3)))};
        \addlegendentry{GELU}
      \end{axis}
    \end{tikzpicture}
  \end{minipage}\hfill
  \begin{minipage}[t]{0.48\textwidth}
    \centering
    \begin{tikzpicture}
      \begin{axis}[
        width=\linewidth, height=5.4cm,
        xlabel={$x \in [0,1]$}, ylabel={$B_{m,3}(x)$},
        xmin=0, xmax=1, ymin=0, ymax=1.12,
        grid=both, grid style={monographRule, opacity=0.4},
        legend pos=north east, legend style={font=\scriptsize, draw=none, fill=white, fill opacity=0.85},
        tick label style={font=\scriptsize}, label style={font=\footnotesize}
      ]
        % Bernstein/Bezier cubic basis on [0,1] (special case of B-spline with 4-fold boundary knots)
        \addplot[domain=0:1, samples=80, monographNavy, thick] {(1-x)^3};
        \addlegendentry{$B_{0,3}$}
        \addplot[domain=0:1, samples=80, monographTeal, thick] {3*x*(1-x)^2};
        \addlegendentry{$B_{1,3}$}
        \addplot[domain=0:1, samples=80, monographIndigo, thick] {3*x^2*(1-x)};
        \addlegendentry{$B_{2,3}$}
        \addplot[domain=0:1, samples=80, monographCrimson, thick] {x^3};
        \addlegendentry{$B_{3,3}$}
        \addplot[domain=0:1, samples=40, black, dashed, thick] {1};
        \addlegendentry{$\sum B_{m,3}$}
      \end{axis}
    \end{tikzpicture}
  \end{minipage}
  \caption{\textbf{Left:} Comparison of ReLU, Softplus, SiLU ($\beta=1$), and GELU activations. \textbf{Right:} Cubic ($k=3$) B-spline basis functions on $[0,1]$ with clamped boundary knots, demonstrating non-negativity and partition of unity ($\sum_{m=0}^3 B_{m,3}(x) = 1$).}
  \label{fig:mlpkan:activations_bsplines}
\end{figure}

Empirical and theoretical region-counting behavior is captured directly in the source literature: Figure~\ref{fig:mlpkan:prince_regions} demonstrates the exponential growth in linear regions per parameter as depth $K$ increases, Figure~\ref{fig:mlpkan:murphy_2d_regions} visualizes the resulting polyhedral tessellation of $\R^2$, and Figure~\ref{fig:mlpkan:kan_arch} illustrates a two-layer KAN with learned univariate edge splines.

\monographfig[0.84\textwidth]{understandingdl_simonprince_p65_fig1.png}{Maximum number of linear regions versus total parameter count for ReLU networks of depth $K \in \{1,\dots,5\}$ with $D_i=1$ and $D_i=10$ inputs, illustrating exponential depth efficiency \bookcite{understandingdl_simonprince}{Fig.~4.7, p.~51}.}{fig:mlpkan:prince_regions}

\monographfig[0.82\textwidth]{probabilisticml_kevinmurphy_introduction_p465_fig1.png}{Decomposition of $\R^2$ into convex linear decision regions produced by a ReLU MLP with (a) one hidden layer of $25$ units and (b) two hidden layers \bookcite{probabilisticml_kevinmurphy_introduction}{Fig.~13.7, p.~435}.}{fig:mlpkan:murphy_2d_regions}

\monographfig[0.68\textwidth]{math_behind_kan_p7_fig1.png}{Two-layer KAN architecture ($[2, 5, 1]$): each edge learns a univariate B-spline activation $\phi_{\ell,q,p}$, while hidden and output nodes perform pure summation \bookcite{math_behind_kan}{PDF p.~7}.}{fig:mlpkan:kan_arch}

% =============================================================================
\section{Numerical Pitfalls, Edge Cases \& Failure Modes}
\label{sec:mlpkan:pitfalls}
% =============================================================================

\begin{pitfall}[Dying ReLU \& Dead-Zone Collapse]
If a ReLU unit's pre-activation satisfies $\vecb{w}_j\T \vecb{x}_n + b_j < 0$ for every training sample $\vecb{x}_n$ (often triggered by an overly large learning rate or negative bias drift), then $\pd{\mathcal{L}}{\vecb{w}_j} = \vecb{0}$ identically and the unit permanently ceases learning.
\textbf{Remedy:} Use He/Kaiming variance-preserving initialization, replace ReLU with LeakyReLU ($\alpha \in [0.01, 0.2]$), GELU, or SiLU, and apply gradient norm clipping.
\end{pitfall}

\begin{pitfall}[Spectral Bias \& Shallow Width Explosion]
Standard ReLU MLPs exhibit a strong inductive bias toward low-frequency functions on low-dimensional coordinate domains, failing to capture fine oscillations without sinusoidal activations or positional encodings \bookcite{probabilisticml_kevinmurphy_introduction}{Sec.~13.2.3, p.~429}. Conversely, relying on shallow UAT widths without compositional depth incurs $O(\epsilon^{-D_i})$ parameter blowup \bookcite{deep_learning_foundations_concepts}{Sec.~6.1.1, p.~173}.
\end{pitfall}

\begin{pitfall}[KAN Boundary Extrapolation \& GPU Throughput Overhead]
Because B-spline basis functions $B_{m,k}(x)$ vanish identically outside the knot support $[t_k, t_{G+k}]$, any out-of-distribution activation collapses $\mathrm{spline}(x) = 0$, leaving only the residual branch $w_b\,\mathrm{SiLU}(x)$ \bookcite{math_behind_kan}{PDF pp.~4--6}.
\textbf{Remedy:} Periodically execute adaptive grid updates (\texttt{update\_grid}) aligned to empirical activation quantiles during training, and reserve KANs for low-dimensional scientific/symbolic tasks where their $5\times$--$10\times$ non-GEMM wall-clock overhead per parameter is offset by compact node counts.
\end{pitfall}

\end{document}
